\chapter*{Introduction Générale}
\markboth{Introduction Générale}{}
\addcontentsline{toc}{chapter}{Introduction Générale}

Une ligne de production n'annonce généralement pas une panne à l'avance. Elle tourne, se dégrade en silence, puis s'arrête, et à ce moment-là le coût ne se limite plus à la réparation : il inclut l'arrêt de production, les contrôles qualité manqués, et dans les pires cas l'exposition à un risque pour la sécurité de la personne qui se trouvait près de la machine. Les systèmes de gestion des actifs industriels bâtis autour de réparations correctives ou d'un calendrier préventif fixe n'ont jamais été conçus pour lire les données de fonctionnement de la machine avant ce point, ce qui est exactement l'écart que ce projet cherche à combler.

C'est cet écart qui a structuré le stage chez SagemCom : une plateforme EAM qui réunit en un seul système la prédiction de panne par apprentissage automatique (Machine Learning), l'accompagnement terrain des techniciens, et un assistant conversationnel connaissant la documentation, plutôt que de les traiter comme trois briques séparées.

La question à laquelle ce rapport répond, concrètement, est la suivante : \emph{une seule plateforme peut-elle porter à la fois le socle opérationnel de la gestion des actifs industriels (utilisateurs, machines, ordres de travail, ordres d'intervention) et une couche prédictive capable d'anticiper les pannes et de guider les équipes de maintenance, sans que l'une ne pénalise l'autre ?}

Le Chapitre~1 expose le contexte du stage, l'entreprise d'accueil et cette problématique en détail. Le Chapitre~2 la traduit en besoins. Le Chapitre~3 conçoit le système qui y répond. Le Chapitre~4 détaille ce qui a été concrètement réalisé. Le Chapitre~5 le met à l'épreuve et rapporte ce qui a tenu et ce qui n'a pas tenu. La conclusion générale revient sur la question ci-dessus et indique clairement dans quelle mesure elle a été résolue.
